% Volume V: Dynamics, Signals, Graphs, and Control
% Continuity note for Chapter 25.
% Previous dependencies: Chapter 8 introduced ODEs, derivatives, differentials, and Euler simulation; Chapter 9 introduced gradients, Jacobians, Hessians, and vector-valued derivatives; Chapter 10 introduced Taylor approximation and local geometry; Chapters 20-24 supplied optimization, numerical methods, matrix exponentials, eigenvalues, and stable computation.
% New durable concepts: initial value problems; autonomous and nonautonomous ODEs; vector fields; flows; existence and uniqueness intuition; equilibria; stability; matrix-exponential solutions; phase portraits; nullclines; nonlinear linearization; Jacobian stability tests; bifurcation intuition; Lyapunov functions; gradient flows; Euler discretization as a bridge between continuous dynamics and iterative algorithms.
% Forward hooks: Chapter 26 extends dynamics from finite-dimensional states to fields; Chapter 27 adds stochastic forcing and SDE notation; Chapters 30 and 36 use dynamical systems for control and reinforcement learning; Chapters 44, 48, and 49 use ODEs for spiking models, latent neural dynamics, attractors, oscillations, and neural fields.
% Notation commitments: States are x(t) in R^n; time derivatives are \dot x(t)=dx/dt; ODEs are \dot x=f(x,t) or autonomous \dot x=f(x); vector fields have Jacobian Df(x) or J_f(x); equilibria are x^* with f(x^*)=0; flows are \Phi_t(x_0); gradients are \nabla E and Hessians are \nabla^2 E; numerical time steps use \Delta t.
\section{Chapter 25: Ordinary Differential Equations and Dynamical Systems}
\label{ch:ordinary-differential-equations-dynamical-systems}

Many systems in computational neuroscience and AI are not single input-output maps. They evolve. A membrane voltage changes over milliseconds, a recurrent network state changes over time, a learning algorithm moves through parameter space, and a latent neural state traces a trajectory during behavior. Ordinary differential equations are the finite-dimensional language for such continuous-time evolution.

The word ordinary means that the unknown state depends on one independent variable, usually time \(t\). The state may be scalar or vector-valued, but it is not a full spatial field. Chapter 26 will add spatial variables and partial derivatives. Chapter 27 will add stochastic forcing. This chapter builds the deterministic baseline.

\paragraph{Chapter problem.}
How do local laws of change \(\dot x=f(x,t)\) determine trajectories, equilibria, stability, phase portraits, and energy-decreasing flows?

\paragraph{Section map.}
Section 25.1 defines ODEs as laws of motion and introduces initial value problems, vector fields, flows, and Euler simulation. Section 25.2 solves and interprets linear dynamical systems. Section 25.3 studies nonlinear dynamics through equilibria, Jacobians, phase portraits, and bifurcations. Section 25.4 develops gradient flows and Lyapunov energy arguments.

\subsection{25.1 ODEs as laws of motion}

\paragraph{Guiding question.}
How can a rule for instantaneous change determine an entire trajectory over time?

\paragraph{Beginner intuition.}
An ordinary differential equation says how fast the state is moving right now. If \(x(t)\) is the state at time \(t\), then
\[
\dot x(t)=\frac{dx}{dt}
\]
is its instantaneous velocity in state space. The ODE
\[
\dot x=f(x,t)
\]
says that the velocity is computed from the current state and possibly the current time. An initial condition
\[
x(t_0)=x_0
\]
selects the particular trajectory.

The smallest example is exponential decay:
\[
\dot x=-x,\qquad x(0)=2.
\]
When \(x>0\), the derivative is negative, so \(x\) decreases. When \(x\) is large, it decreases quickly. The solution is
\[
x(t)=2e^{-t}.
\]
The ODE gives the local law; the solution stitches those local velocities into a curve.

\paragraph{Formal setup.}
Let \(D\subseteq\mathbb{R}^n\) be a state domain and let \(I\subseteq\mathbb{R}\) be a time interval. A first-order ODE is
\[
\dot x(t)=f(x(t),t),
\qquad
f:D\times I\to\mathbb{R}^n.
\]
An initial value problem is
\[
\dot x(t)=f(x(t),t),\qquad x(t_0)=x_0.
\]
If \(f\) does not explicitly depend on \(t\), the system is autonomous:
\[
\dot x=f(x).
\]
For autonomous systems, \(f\) is called a vector field because it assigns a velocity vector \(f(x)\) to each state \(x\).

When a unique solution exists, the flow map is written
\[
\Phi_t(x_0)=x(t),
\]
where \(x(t)\) is the solution starting at \(x(0)=x_0\). For autonomous systems, flows satisfy
\[
\Phi_{t+s}(x_0)=\Phi_t(\Phi_s(x_0))
\]
whenever both sides are defined.

A standard existence and uniqueness result says: if \(f(x,t)\) is continuous in \(t\) and locally Lipschitz in \(x\), then for each initial condition there is a unique local solution. This is a theorem about deterministic predictability in the mathematical model, not a claim that biological systems are noise-free.

\paragraph{Worked mechanism: Euler simulation from differentials.}
The notation
\[
\dot x(t)=f(x(t),t)
\]
can be read as
\[
dx\approx f(x,t)\,dt
\]
over a small time interval, with the understanding that this is ordinary deterministic calculus. Replacing \(dt\) by a finite step \(\Delta t\) gives forward Euler:
\[
x_{k+1}=x_k+\Delta t\,f(x_k,t_k).
\]
For \(\dot x=-x\), \(x_0=2\), and \(\Delta t=0.1\),
\[
x_1=2+0.1(-2)=1.8,\qquad
x_2=1.8+0.1(-1.8)=1.62.
\]
Euler approximates the exact values \(2e^{-0.1}\) and \(2e^{-0.2}\). It turns the continuous law into a computable recurrence.

\paragraph{Examples and connections.}
A numerical symbolic example is
\[
\dot x=2t,\qquad x(0)=1.
\]
Integrating gives
\[
x(t)=1+t^2.
\]
Here the velocity depends on time rather than state, so the system is nonautonomous.

In computational neuroscience, the passive membrane equation
\[
\tau \dot V=-(V-E_L)+RI(t)
\]
describes voltage relaxation toward a current-dependent steady value. The state is \(V(t)\), the derivative \(\dot V\) is measured in voltage per unit time, and \(\tau\) sets the response timescale.

In AI, a residual network layer
\[
h_{k+1}=h_k+\Delta t\,f_\theta(h_k)
\]
can be read as an Euler step for a continuous-depth model
\[
\dot h=f_\theta(h).
\]
This does not mean every residual network should be interpreted as a literal physical flow, but the mathematical analogy is useful for stability and solver design.

\paragraph{Common misconceptions and failure modes.}
The symbol \(dx/dt\) is not a fraction of two ordinary numbers, but in deterministic calculus it records a limiting rate and can guide valid differential reasoning when smoothness assumptions hold. An ODE does not determine a unique trajectory without an initial condition. A slope field is not a solution; it is the velocity rule from which solution curves are drawn. Finally, numerical simulation is not the same as the exact flow. Step size and method matter.

\paragraph{Exercises.}
\begin{enumerate}[label=\textbf{25.1.\arabic*.}, leftmargin=*]
\item \textbf{Mechanical.} For \(\dot x=3x\), \(x(0)=2\), write the exact solution.
\item \textbf{Proof or derivation.} Show that if \(x(t)=x_0e^{at}\), then \(x\) solves \(\dot x=ax\) with \(x(0)=x_0\).
\item \textbf{Numerical/computational.} Use two Euler steps with \(\Delta t=0.2\) for \(\dot x=-2x\), \(x_0=1\). Compare with \(e^{-0.8}\) at \(t=0.4\).
\item \textbf{Interpretation.} In the membrane equation, what does a larger \(\tau\) mean for the speed of voltage relaxation?
\item \textbf{Misconception/debugging.} A student writes an ODE but no initial condition and claims to have a unique trajectory. Explain what is missing.
\end{enumerate}

\subsection{25.2 Linear dynamical systems}

\paragraph{Guiding question.}
How do matrices determine the trajectories, modes, and stability of linear continuous-time systems?

\paragraph{Beginner intuition.}
A linear dynamical system has a velocity proportional to the state:
\[
\dot x=Ax.
\]
The matrix \(A\) tells which directions grow, decay, rotate, or mix. This is the dynamic version of the eigenvector story from Chapter 6. If \(Av=\lambda v\), then along direction \(v\), the system behaves like the scalar equation \(\dot c=\lambda c\). Real negative eigenvalues decay, real positive eigenvalues grow, and complex eigenvalues create rotations with growth or decay.

\paragraph{Formal setup.}
For \(A\in\mathbb{R}^{n\times n}\), the homogeneous linear system is
\[
\dot x(t)=Ax(t),\qquad x(0)=x_0.
\]
Its solution is
\[
\boxed{x(t)=e^{At}x_0,}
\]
where the matrix exponential is
\[
e^{At}=I+At+\frac{(At)^2}{2!}+\frac{(At)^3}{3!}+\cdots.
\]
If \(A=V\Lambda V^{-1}\), then
\[
e^{At}=Ve^{\Lambda t}V^{-1},
\]
where \(e^{\Lambda t}\) is diagonal with entries \(e^{\lambda_i t}\).

An input-driven linear system has the form
\[
\dot x=Ax+Bu(t).
\]
The solution is
\[
x(t)=e^{A(t-t_0)}x(t_0)+
\int_{t_0}^t e^{A(t-s)}B u(s)\,ds.
\]
The first term propagates the initial condition; the integral accumulates the effect of input over time.

\paragraph{Worked mechanism: stability from eigenvalues.}
For the homogeneous system \(\dot x=Ax\), the equilibrium \(x^*=0\) is asymptotically stable if all eigenvalues of \(A\) have negative real parts. It is unstable if at least one eigenvalue has positive real part. If eigenvalues lie on the imaginary axis, linear stability requires more care.

Example:
\[
A=\begin{bmatrix}-1&0\\0&-3\end{bmatrix}.
\]
Then
\[
e^{At}=\begin{bmatrix}e^{-t}&0\\0&e^{-3t}\end{bmatrix},
\qquad
x(t)=\begin{bmatrix}e^{-t}x_1(0)\\e^{-3t}x_2(0)\end{bmatrix}.
\]
The second coordinate decays faster.

For
\[
A=\begin{bmatrix}0&-\omega\\ \omega&0\end{bmatrix},
\]
the eigenvalues are \(\pm i\omega\), and trajectories rotate without decay in the ideal linear model. Adding \(-\alpha I\) with \(\alpha>0\) gives a stable spiral.

\paragraph{Examples and connections.}
A two-dimensional numerical system
\[
\dot x=
\begin{bmatrix}-1&2\\0&-2\end{bmatrix}x
\]
has eigenvalues \(-1\) and \(-2\), so all trajectories decay to the origin, although the off-diagonal term shears the phase portrait.

In computational neuroscience, a linear rate model near a baseline can be written
\[
\tau \dot r=-r+Wr+I_{\mathrm{ext}},
\]
where \(r\in\mathbb{R}^n\) is the activity vector,
\(I_{\mathrm{ext}}\in\mathbb{R}^n\) is an external input, and
\(I_n\in\mathbb{R}^{n\times n}\) is the identity matrix. Equivalently,
\[
\dot r=\frac{1}{\tau}(W-I_n)r+\frac{1}{\tau}I_{\mathrm{ext}}.
\]
The eigenvalues of \(\tau^{-1}(W-I_n)\) determine whether small activity perturbations decay, grow, or oscillate.

In AI, linear recurrent networks and state-space models use equations such as
\[
h_{t+1}=Ah_t+Bu_t
\]
in discrete time or \(\dot h=Ah+Bu(t)\) in continuous time. Stability controls whether hidden states remember, forget, or explode.

\paragraph{Common misconceptions and failure modes.}
Eigenvalues alone are not always the full numerical story for nonnormal matrices; transient amplification can occur even when all eigenvalues are stable. A second misconception is that complex eigenvalues mean complex-valued states. Real matrices with complex-conjugate eigenvalues generate real oscillatory trajectories. A third failure is to apply the matrix exponential formula without checking that the system is time-invariant; if \(A=A(t)\), the solution is more complicated.

\paragraph{Exercises.}
\begin{enumerate}[label=\textbf{25.2.\arabic*.}, leftmargin=*]
\item \textbf{Mechanical.} Solve \(\dot x=Ax\) for \(A=\operatorname{diag}(-2,1)\) and \(x(0)=(3,4)^\top\).
\item \textbf{Proof or derivation.} Show from the power series that \(\frac{d}{dt}e^{At}=Ae^{At}\) for constant \(A\).
\item \textbf{Numerical/computational.} Classify the origin for systems with eigenvalues \((-1,-4)\), \((2,-1)\), and \((-0.5\pm 3i)\).
\item \textbf{Interpretation.} In a linearized recurrent neural population model, what does an eigenvalue with real part close to zero imply about timescale?
\item \textbf{Misconception/debugging.} A student says, "A stable linear system can never show transient growth." Explain why nonnormal matrices complicate this statement.
\end{enumerate}

\subsection{25.3 Nonlinear dynamics}

\paragraph{Guiding question.}
How can we analyze systems whose velocity depends nonlinearly on the current state?

\paragraph{Beginner intuition.}
Nonlinear systems can have multiple equilibria, thresholds, oscillations, and bifurcations. Exact solutions are usually unavailable. Instead, we study the geometry of the vector field: where it vanishes, which way arrows point, how trajectories move near equilibria, and how the behavior changes when parameters vary.

The local tool is linearization. Near a state \(x^*\), Taylor expansion gives
\[
f(x)\approx f(x^*)+Df(x^*)(x-x^*).
\]
If \(x^*\) is an equilibrium, then \(f(x^*)=0\), and nearby perturbations approximately follow the linear system
\[
\dot \xi=Df(x^*)\xi,
\qquad \xi=x-x^*.
\]

\paragraph{Formal setup.}
An autonomous nonlinear system is
\[
\dot x=f(x),\qquad f:D\subseteq\mathbb{R}^n\to\mathbb{R}^n.
\]
An equilibrium is a point \(x^*\in D\) such that
\[
f(x^*)=0.
\]
The Jacobian at \(x^*\) is
\[
J_f(x^*)=Df(x^*)\in\mathbb{R}^{n\times n},
\qquad
(J_f)_{ij}=\frac{\partial f_i}{\partial x_j}(x^*).
\]
If all eigenvalues of \(J_f(x^*)\) have negative real parts, then \(x^*\) is locally asymptotically stable. If at least one eigenvalue has positive real part, then \(x^*\) is unstable. This is a local theorem; it describes behavior near the equilibrium, not necessarily far away.

For a two-dimensional system
\[
\dot x=f(x,y),\qquad \dot y=g(x,y),
\]
the \(x\)-nullcline is \(f(x,y)=0\), and the \(y\)-nullcline is \(g(x,y)=0\). Equilibria occur at intersections of nullclines. Phase portraits combine nullclines, arrows, and representative trajectories.

\paragraph{Worked mechanism: logistic growth and bifurcation intuition.}
The logistic equation is
\[
\dot x=rx\left(1-\frac{x}{K}\right),
\qquad r>0,\quad K>0.
\]
Equilibria satisfy
\[
rx\left(1-\frac{x}{K}\right)=0,
\]
so \(x^*=0\) and \(x^*=K\). The derivative of the vector field is
\[
f'(x)=r-\frac{2r}{K}x.
\]
At \(x^*=0\), \(f'(0)=r>0\), so the equilibrium is unstable. At \(x^*=K\), \(f'(K)=-r<0\), so it is locally stable.

A simple bifurcation example is
\[
\dot x=\mu-x^2.
\]
Equilibria satisfy \(x=\pm\sqrt{\mu}\) when \(\mu>0\). For \(\mu<0\), there are no real equilibria. At \(\mu=0\), two equilibria collide. This is a saddle-node bifurcation. The lesson is that changing a parameter can change the qualitative structure of the phase portrait.

\paragraph{Examples and connections.}
A symbolic two-dimensional system
\[
\dot x=x(1-x)-y,\qquad
\dot y=-y+x
\]
has nullclines \(y=x(1-x)\) and \(y=x\). Their intersections are candidate equilibria. Linearizing at each intersection gives local stability information.

In computational neuroscience, Wilson-Cowan rate models use nonlinear activation functions:
\[
\tau_E\dot E=-E+S(w_{EE}E-w_{EI}I+P),
\]
\[
\tau_I\dot I=-I+S(w_{IE}E-w_{II}I+Q).
\]
Nullclines organize the excitatory-inhibitory phase plane. Depending on parameters, the system can settle to a fixed point, show switching, or support oscillatory tendencies. The equations are models, not direct claims that two variables capture all cortical dynamics.

In AI, nonlinear recurrent networks
\[
\dot h=-h+\sigma(Wh+Ux)
\]
can have multiple attractors or long transients. Linearization helps diagnose whether a learned latent state is stable under small perturbations.

\paragraph{Common misconceptions and failure modes.}
Linearization is local. It cannot prove global behavior by itself. A stable Jacobian near one equilibrium does not rule out another attractor elsewhere. A second misconception is that equilibria are the only important behavior; limit cycles, chaotic attractors, and transient dynamics may matter. A third failure is to draw nullclines and forget arrows; nullclines show where one component stops changing, not the full direction of motion.

\paragraph{Exercises.}
\begin{enumerate}[label=\textbf{25.3.\arabic*.}, leftmargin=*]
\item \textbf{Mechanical.} Find the equilibria of \(\dot x=x(2-x)\) and classify them using \(f'(x)\).
\item \textbf{Proof or derivation.} Starting from Taylor expansion, derive the linearized system \(\dot \xi=Df(x^*)\xi\) near an equilibrium.
\item \textbf{Numerical/computational.} For \(\dot x=y\), \(\dot y=-x-0.2y\), compute the Jacobian and its eigenvalues. Classify the origin.
\item \textbf{Interpretation.} In a Wilson-Cowan phase plane, what does an intersection of the \(E\)- and \(I\)-nullclines represent?
\item \textbf{Misconception/debugging.} A student claims a nonlinear system is globally stable because the Jacobian at one equilibrium has negative eigenvalues. Explain the missing global argument.
\end{enumerate}

\subsection{25.4 Gradient flows and energy landscapes}

\paragraph{Guiding question.}
What happens when a dynamical system always moves downhill in an energy or loss landscape?

\paragraph{Beginner intuition.}
A gradient flow is the continuous-time version of steepest descent. Given a differentiable energy or loss
\[
E:\mathbb{R}^n\to\mathbb{R},
\]
the gradient flow is
\[
\dot x=-\nabla E(x).
\]
The state moves in the direction of fastest local decrease of \(E\). This creates a direct bridge between dynamics and optimization.

\paragraph{Formal setup.}
Let \(E\) be continuously differentiable. A trajectory \(x(t)\) follows the gradient flow if
\[
\frac{dx}{dt}=-\nabla E(x(t)).
\]
Along any smooth trajectory satisfying this equation,
\[
\frac{d}{dt}E(x(t))
=\nabla E(x(t))^\top \dot x(t)
=-\|\nabla E(x(t))\|_2^2
\leq 0.
\]
Thus \(E\) is nonincreasing along trajectories. If \(\nabla E(x)\neq 0\), the energy strictly decreases at that instant.

A Lyapunov function for a dynamical system is a scalar function that decreases along trajectories and helps prove stability. For gradient flow, the energy itself is a natural Lyapunov function.

\paragraph{Worked mechanism: quadratic energy.}
Let
\[
E(x)=\frac12 x^\top Hx,
\]
where \(H=H^\top\) is positive definite. Then
\[
\nabla E(x)=Hx,
\]
and the gradient flow is
\[
\dot x=-Hx.
\]
If \(H=V\Lambda V^\top\), then in eigen-coordinates \(z=V^\top x\),
\[
\dot z_i=-\lambda_i z_i,
\qquad
z_i(t)=z_i(0)e^{-\lambda_i t}.
\]
Large eigenvalues create fast decay; small eigenvalues create slow directions. This is the continuous-time version of conditioning in optimization.

Forward Euler applied to the gradient flow gives
\[
x_{k+1}=x_k-\Delta t\,\nabla E(x_k),
\]
which is ordinary gradient descent. The step size \(\Delta t\) is a numerical time step. If it is too large, the discrete iteration can overshoot even though the continuous gradient flow decreases \(E\) monotonically.

\paragraph{Examples and connections.}
A one-dimensional energy
\[
E(x)=\frac14x^4-\frac12x^2
\]
has gradient
\[
E'(x)=x^3-x.
\]
The gradient flow is
\[
\dot x=-(x^3-x)=x-x^3.
\]
Equilibria are \(x=-1,0,1\). The points \(-1\) and \(1\) are stable minima, while \(0\) is an unstable maximum. The basins of attraction are separated by the unstable equilibrium.

In computational neuroscience, Hopfield-style models use an energy function to describe attractor memory in an idealized network. The precise theorem depends on symmetric weights and asynchronous updates or appropriate continuous dynamics. The safe statement is: under those modeling assumptions, the energy decreases and stable patterns correspond to attractors.

In AI, optimization dynamics can be approximated by gradient flow when the step size is small. Loss landscapes may contain basins, saddles, flat directions, and sharp directions. The gradient-flow calculation explains why a loss decreases continuously, but discrete stochastic training adds step-size effects and gradient noise.

\paragraph{Common misconceptions and failure modes.}
"Energy" is a mathematical Lyapunov function here, not necessarily physical energy. A gradient flow cannot cycle in a region where \(E\) strictly decreases, but non-gradient systems can oscillate. A local minimum is stable for gradient flow under suitable curvature assumptions, but it may not be globally optimal. Finally, gradient descent is not identical to gradient flow; it is a time discretization with its own stability constraints.

\paragraph{Exercises.}
\begin{enumerate}[label=\textbf{25.4.\arabic*.}, leftmargin=*]
\item \textbf{Mechanical.} For \(E(x,y)=\frac12(x^2+4y^2)\), write the gradient flow equations.
\item \textbf{Proof or derivation.} Prove that \(dE(x(t))/dt=-\|\nabla E(x(t))\|^2\) along \(\dot x=-\nabla E(x)\).
\item \textbf{Numerical/computational.} Apply one gradient-descent step with step size \(0.1\) to \(E(x)=\frac12(5x^2)\) starting at \(x_0=2\). Compare with one Euler step for the gradient flow.
\item \textbf{Interpretation.} In an attractor memory model with an energy function, what does a basin of attraction represent?
\item \textbf{Misconception/debugging.} A student says, "Because gradient flow decreases energy, gradient descent with any step size must decrease energy." Give a reason this is false.
\end{enumerate}

\paragraph{Closing bridge.}
ODEs describe finite-dimensional deterministic motion. They give the baseline language for trajectories, stability, attractors, and numerical simulation. Chapter 26 extends the state from a finite vector \(x(t)\) to a field \(u(x,t)\) over space and time. Chapter 27 then changes the differential notation again by adding stochastic increments such as \(dW_t\), where ordinary deterministic calculus no longer applies without new rules.

\clearpage
